\bookchapter{第三十一章：短影片拍攝與內容創作入門}

\chapterintro{2026年，短影片已經不是“年輕人的玩具”了，它變成了所有人獲取資訊和表達自我的基礎方式。不管你是想記錄生活、分享知識，還是做點副業，手機就是你最好的創作工具——不需要買相機、不需要學Premiere，一部手機就能開始。這一章幫你從零起步。}

\booksection{一、拍之前先想清楚：你要拍給誰看？}

很多人拍影片的第一個錯誤就是“拿起手機就拍”，拍完不知道給誰看。在開拍之前，先回答三個問題：

\begin{enumerate}
  \item \textbf{我的觀眾是誰？}——拍給孩子看、拍給同齡人看、拍給長輩看，內容風格完全不同。
  \item \textbf{我能提供什麼價值？}——讓人笑（娛樂）、讓人學（知識）、讓人饞（美食）、讓人暖（情感）。至少佔一樣。
  \item \textbf{我能在哪個平臺持續更新？}——抖音/快手（泛娛樂）、小紅書（生活方式）、B站（深度內容）、影片號（社交圈）。選一個主陣地深耕，不要一上來就全平臺分發。
\end{enumerate}

\booksection{二、手機拍攝基礎設定——開拍前必做的5件事}

\begin{enumerate}
  \item \textbf{解析度設1080P，幀率設30fps}：這是短影片平臺的標準。4K檔案太大上傳慢，60fps在社交平臺上沒有優勢反而檔案更大。
  \item \textbf{擦乾淨鏡頭}：第一條也是最重要的一條。拍之前用衣服擦一下鏡頭，畫面立刻清晰一個檔次。
  \item \textbf{鎖定對焦和曝光}：在取景框長按畫面主體2秒，會出現"AE/AF鎖定"提示。這樣拍的時候不會因為移動導致畫面忽明忽暗。
  \item \textbf{橫拍還是豎拍？}——發抖音/快手/影片號/小紅書\textbf{豎著拍}（9:16），發B站/YouTube\textbf{橫著拍}（16:9）。想清楚平臺再決定方向。
  \item \textbf{開網格線}：設定 > 相機 > 網格線。三分法構圖的基礎，幫你快速找到畫面重心。
\end{enumerate}

\booksection{三、運鏡五招——讓畫面“動”起來}

新手拍影片最大的問題是畫面太“靜”——全程站著不動拍。學會下面五個基本運鏡，畫面立刻有專業感：

\begin{itemize}
  \item \textbf{推（Push）}：手機慢慢向被拍物件靠近。製造"靠近"的緊張感或親密感。
  \item \textbf{拉（Pull）}：手機慢慢遠離被拍物件。從細節拉到全景，交代環境。
  \item \textbf{搖（Pan）}：手機位置不動，左右轉動。適合拍風景、展示空間。
  \item \textbf{移（Track）}：手機平行移動，跟著被拍物件走。需要手持穩定，最好配合手機穩定器。
  \item \textbf{跟（Follow）}：手機跟在被拍物件後面或前面走。Vlog最常用的運鏡。
\end{itemize}

\textbf{穩是第一要務}：手抖的畫面看起來很不專業。雙手持機比單手穩；走路時膝蓋微屈，用“忍者步”減少上下顛簸；如果有預算，買一個手機穩定器（大疆OM系列約500-800元），效果質變。

\booksection{四、剪輯入門——剪映就夠了}

剪映（CapCut的國內版）是目前最適合新手的手機剪輯工具，免費、功能全、模板多。以下是基礎工作流：

\begin{enumerate}
  \item \textbf{匯入素材}：開啟剪映 > 開始創作 > 選擇拍好的影片片段匯入。
  \item \textbf{裁剪}：拖動時間軸，把每個片段開頭結尾的廢片剪掉。原則：每一段只保留有資訊量的部分，寧短勿長。
  \item \textbf{加轉場}：兩段影片之間點"轉場"按鈕，選一個簡單的轉場效果。新手建議用"疊化"或"黑場"，不要用花哨的3D轉場。
  \item \textbf{加字幕}：剪映有"識別字幕"功能，一鍵自動生成字幕。準確率約90\%，手動校對一遍即可。
  \item \textbf{配音樂}：剪映音樂庫裡有大量版權免費的音樂。選一首節奏和內容匹配的BGM，音量調到影片原聲的1/3左右。
  \item \textbf{匯出}：1080P匯出，幀率30fps，位元速率選"推薦"即可。
\end{enumerate}

\booksection{五、封面的重要性——70\%的人只看封面決定看不看}

短影片的封面就是“門面”，決定別人要不要點進來。幾個做封面的要點：

\begin{itemize}
  \item \textbf{封面要有一句話}：用大字寫明影片的核心內容，比如"3分鐘學會手機修圖"。字型要大、顏色要亮、和背景有對比。
  \item \textbf{封面和內容要一致}：不要做"標題黨"——封面說一件事，點進去是另一回事，觀眾會立刻劃走，演演算法會降權。
  \item \textbf{人物表情要生動}：如果封面有人，選一個表情生動或誇張的瞬間。人臉比文字更能吸引點選。
  \item \textbf{剪映做封面}：在剪映裡選一幀最好的畫面，加文字疊加，匯出為圖片作為封面。
\end{itemize}

\booksection{六、內容創作的三條心法}

\begin{itemize}
  \item \textbf{前3秒定生死}：短影片平臺，觀眾3秒內決定劃走還是繼續看。前3秒必須給出"鉤子"——一個問題、一個懸念、一個意外畫面。
  \item \textbf{一條影片只講一件事}：新手最容易犯的錯是一條影片塞太多內容。一條影片只解決一個問題，講透比講全重要。
  \item \textbf{堅持比天賦重要}：演演算法推薦需要"養號"。前20條影片幾乎不會有播放量，這是正常的。堅持每週更新2-3條，讓演演算法認識你的內容方向，流量會慢慢起來。
\end{itemize}

\begin{tcbpractice}{動手：拍一條30秒的短影片}
\begin{enumerate}
  \item 選一個你擅長的小知識（比如"怎麼把手機字型調大"）。
  \item 用手機豎屏拍攝3-5段素材，每段5-10秒。
  \item 用剪映拼接、裁剪、加字幕、配BGM。
  \item 做一張帶大字的封面圖。
  \item 匯出後先發給朋友看，收集反饋再發布。
\end{enumerate}

\end{tcbpractice}

\begin{tcbsummary}{本章小結}
\begin{itemize}
  \item 拍之前想清楚：觀眾是誰、提供什麼價值、在哪個平臺發
  \item 開拍5件事：1080P/30fps、擦鏡頭、鎖對焦曝光、確定橫豎、開網格線
  \item 運鏡五招：推、拉、搖、移、跟——穩是第一要務
  \item 剪輯用剪映：裁剪轉場字幕BGM匯出，六步走
  \item 前3秒定生死，一條影片只講一件事，堅持比天賦重要
\end{itemize}

\end{tcbsummary}
